\section{Kondisi dan Fungsi Kerja}

% \begin{figure}[h]
%     \centering
%     
\includegraphics[width=7.5cm]{dummy.png}
%     \caption{Lorong Ruang Kerja}
%     \label{fig:lorongruangkerja}
% \end{figure}

% \begin{figure}[h]
%     \centering
%     
\includegraphics[width=7.5cm]{dummy.png}
%     \caption{Ruang Pantry}
%     \label{fig:ruangpantry}
% \end{figure}

% \begin{figure}[h]
%     \centering
%     
\includegraphics[width=7.5cm]{dummy.png}
%     \caption{Ruang Rapat}
%     \label{fig:ruangrapat}
% \end{figure}

% \begin{figure}[h]
%     \centering
%     
\includegraphics[width=7.5cm]{dummy.png}
%     \caption{Ruang Auditorium}
%     \label{fig:ruangauditorium}
% \end{figure}

Pelaksanaan magang industri di PT Inti Utama Solusindo (PT IUS) memberikan
eksposur langsung terhadap ekosistem pengembangan perangkat lunak skala
korporasi yang terstruktur. Kondisi kerja di perusahaan ini didesain untuk
menyeimbangkan kedisiplinan operasional dengan fleksibilitas kreatif yang
dibutuhkan oleh seorang pengembang.

\subsection{Penempatan dan Ruang Kerja}
Kegiatan magang dilaksanakan pada Departemen Hospinet / Pharmalogic PT Inti
Utama Solusindo yang berlokasi di Lantai 2, Gedung Century (Pharos Group).
Departemen ini mengusung konsep \textit{open-space office} yang bertujuan untuk
mengurangi hambatan komunikasi antar tim \textit{developer}, QA, dan PM./@

\begin{figure}[H]
      \centering
      \includegraphics[width=7.5cm]{camera-ruangkerja.png}
      \caption{Ruang Kerja}%
      \label{fig:ruangkerja}
\end{figure}

\subsection{Fungsi dan Peran}
Peran yang dijalankan selama magang adalah \textit{Front-End Developer}, dengan
tanggung jawab utama mengimplementasikan tampilan antarmuka berdasarkan desain
yang sudah diberikan oleh Tim UI/UX.\@ Pekerjaan sehari-hari mencakup pembuatan
layar-layar baru menggunakan React, peningkatan fitur \textit{widget/component}
dan perbaikan \textit{bug} pada fitur-fitur yang sudah berjalan.

Mentor sekaligus \textit{Front-End Developer Supervisor} memiliki tanggung
jawab untuk mengawasi pembuatan kode yang dibuat oleh tim FE.\@ Interaksi
dengan mentor dilakukan pada saat dibutuhkan tindakan yang membutuhkan peran
\textit{supervisor}, seperti \textit{code review, merge request,} dan
menanyakan kode yang sulit dimengerti dalam \textit{framework}. Untuk pemberian
tugas (\textit{Task}) dan pemantauan progres, koordinasi dilakukan bersama
\textit{Deputy Project Manager} (DPM) melalui platform OpenProject.

Dalam pengerjaan proyek yang terkait pembuatan layar baru, seringkali dilakukan
koordinasi dengan tim \textit{Back-End} untuk integrasi data dan API.\@ Apabila
terkait pengerjaan perbaikan \textit{bug} atau \textit{improvement} fitur baru,
dilakukan koordinasi rutin dengan tim \textit{Quality Assurance} untuk
pengujian fitur.

\subsection{Sistem Presensi}
Sistem pencatatan kehadiran di PT IUS menggunakan mesin fingerprint yang
tersedia di kantor. Setiap karyawan dan intern diwajibkan melakukan rekam sidik
jari saat tiba (finger-in) dan saat akan pulang (finger-out). Jam kerja resmi
dimulai pukul 08.00 WIB, dengan toleransi keterlambatan hingga pukul 10.00 WIB.

\begin{figure}[H]
      \centering
      \includegraphics[width=7.5cm]{fingerprint.png}
      \caption{Mesin Fingerprint}%
      \label{fig:mesinfingerprint}
\end{figure}

Apabila penulis tiba setelah batas waktu tersebut, kehadiran pada hari itu
tidak akan tercatat sebagai hari masuk oleh sistem. Selain aturan jam masuk,
terdapat juga ketentuan minimum durasi kerja harian sebesar sembilan jam.
Artinya, meskipun hadir tepat waktu, apabila total durasi kerja dalam satu hari
kurang dari sembilan jam — misalnya masuk pukul 08.30 lalu pulang pukul 17.00 —
sistem secara otomatis tidak menghitung hari tersebut sebagai hari kerja penuh.
Ketentuan ini mendorong kedisiplinan tidak hanya dari sisi waktu kedatangan,
tetapi juga dari sisi konsistensi durasi kerja.

\begin{figure}[H]
      \centering
      \includegraphics[width=7.5cm]{screen-sipp-main.png}
      \caption{Dashboard SIPP Karyawan Pharos}
\end{figure}

Rekap kehadiran dapat dipantau secara mandiri melalui portal Sistem Informasi
Pelaporan Perusahaan (SIPP) di https://sipp.pharmalink.id/, yang menampilkan
informasi waktu masuk, waktu keluar, durasi harian, dan status kehadiran secara
keseluruhan.

\begin{figure}[H]
      \centering
      \includegraphics[width=7.5cm]{screen-sipp-kehadiran.png}
      \caption{Daftar Kehadiran Pribadi}
\end{figure}

\subsection{Budaya Berpakaian}
Perusahaan memberlakukan aturan berpakaian yang cukup standar untuk lingkungan
perusahaan pada umumnya. Senin hingga Kamis menggunakan pakaian kemeja dan
celana bahan (formal). Hari Jumat diberlakukan sebagai \textit{Casual Friday},
di mana karyawan boleh mengenakan pakaian yang lebih santai seperti kaos polo
dan \textit{jeans}.

Pharos juga memiliki peraturan unik dimana setiap karyawan wajib melepaskan
sepatu yang dikenakan dari luar lalu menggantinya dengan sepatu khusus dalam
kantor (berupa sepatu Crocs).

\subsection{Konsumsi Siang Gratis}
Seluruh karyawan dan peserta magang mendapatkan makan siang gratis dari kantin
Pharos Group. Secara tidak langsung dapat menjadi kesempatan untuk berinteraksi
dengan rekan-rekan dari tim lain di ruang makan.

\subsection{Fasilitas Lainnya}
Beberapa fasilitas pendukung lain yang tersedia di kantor:
\begin{enumerate}[noitemsep, topsep=2pt]
      \item WiFi: Jaringan internet tersedia melalui dua jalur frekuensi, `Pharmalogic
            2.4GHz' dan `Pharmalogic 5GHz', yang cukup stabil untuk kebutuhan
            \textit{coding, deployment}, maupun untuk \textit{download package} berukuran
            besar.
      \item Ruang Pertemuan: Tersedia ruang rapat untuk koordinasi tim dalam skala kecil,
            serta Ruang Auditorium untuk pertemuan yang melibatkan lebih banyak peserta,
            seperti presentasi hasil pekerjaan divisi ke \textit{supervisor}.
      \item Ruang Auditorium: Di Lantai 4 Gedung Century terdapat ruang pertemuan besar
            yang digunakan pada saat ada koordinasi tim yang memerlukan semua anggota
            perusahaan ikut serta, seperti \textit{Objective Alignment} dengan direksi.
      \item Ruang Pantry: Tersedia ruang untuk mencuci peralatan makan/minum dan membuang
            sampah basah.
\end{enumerate}

Secara keseluruhan, kondisi kerja di PT IUS didesain untuk menciptakan
lingkungan yang sustainable, di mana fasilitas fisik yang memadai bersinergi
dengan budaya kerja yang disiplin untuk menghasilkan produk digital berkualitas
tinggi.